\section{관련 연구 및 문제 정의}

\subsection{회전 다리와 바퀴 다리}

RHex는 각 힙에 연속 회전하는 탄성 다리 하나만 두고, 시계 기반 삼각보행으로
기계적 단순성과 험지 이동성을 함께 달성했다
\cite{saranli2001rhex,altendorfer2001rhex}. 이후 전용 제어기를 통해 일반 크기의
계단을 반복 등반하는 성능도 입증되었다 \cite{moore2002reliable}. X-RHex 계열은
이 구조를 범용 연구 플랫폼으로 확장했다 \cite{galloway2010xrhex}. Whegs는 탄성 다중 스포크와 결합
구동을 이용해 구동기 수를 줄였으며 \cite{quinn2002whegs}, Q-Whex는 준바퀴와
계단 형상에 맞춘 위상 패턴을 사용했다 \cite{zhang2023qwhex}. C자 다리를 사용하는
6족 로봇의 운동 및 동력 요구량도 모델링 후 실물 시험대에서 검증되었다
\cite{vina2021cleg}. RHex의 반원형
다리는 본 연구의 말단 형상과 가장 유사하므로 두 구조의 차이를 명확히 구분해야
한다. RHex의 다리는 하나의 힙 회전으로 구동되는 탄성 C자 구조이기 때문에
접촉점의 이동과 몸체 자세가 동일한 입력에 의해 결정된다. 반면 본 연구의 C자
호는 3자유도 다리의 세 번째 관절이다. 따라서 호를 멈추거나 원하는 위치에
배치하거나 회전시키는 동안에도 나머지 관절로 몸체 자세를 조절할 수 있다. 절제
실험에서 검증하려는 핵심은 C자 형상 자체가 아니라 이러한 제어 권한의 분리다.

두 번째 연구군은 기구의 형태를 직접 바꾼다. Wheel Transformer는 장애물과
접촉하면 수동 스포크를 펼치고 \cite{kim2014wheeltransformer}, Quattroped와
TurboQuad는 전용 변형 자유도를 이용해 다리를 바퀴로 전환한다
\cite{chen2014quattroped,chen2017turboquad}. 두 이동 방식을 모두 확보할 수
있지만, 별도의 변형 기구와 구동계가 필요하다. 세 번째 연구군은 관절 다리 끝에
구동 바퀴를 추가한다. ASTERISK H는 감지한 장애물에 따라 구름 이동과 삼각보행을
연속적으로 전환한다 \cite{yoshioka2008asterisk}. Cassino Hexapod III는 인체형 다리 끝에 옴니휠을 장착했고
\cite{tedeschi2017cassino}, ANYmal on Wheels는 바퀴의 구름 구속을 전신 최적화에
포함했다 \cite{bjelonic2019keep,bjelonic2021wholebody,hutter2016anymal}.
이륜 점프 로봇도 작은 구조 안에서 다리와 바퀴의 제어 권한을 분리한다
\cite{klemm2019ascento}. 이러한 연구는 다리와 바퀴를 동시에 제어할 때의 이점을
보여주지만, 다리 관절과 별도로 바퀴 구동기를 유지한다. 본 연구는 하이브리드
이동 자체의 신규성을 주장하지 않는다. 검증하려는 차별점은 \emph{액추에이터와
접촉부의 재사용}이다. 하나의 말단 관절과 개방형 호가 발 궤적, 구름 운동, 계단
모서리 걸림을 모두 만들며, 별도의 변형 상태나 추가 구동기가 필요하지 않다.

\subsection{운용 범위의 정식화}

몸체 명령과 모드를
\begin{equation}
 \vect{c}=[v_x,\;v_y,\;\omega_z]^\top\in\mathcal{C},\qquad
 m\in\{W,R,C\}                                                 \label{eq:command-ko}
\end{equation}
로 정의한다. 여기서 $\vect{q}\in\R^{18}$은 관절 벡터이고, $\Gamma_i$는 $i$번째
호에서 활성화된 접촉의 집합이다. 모드에 따라 기구 $\mathcal{M}_m$을 바꾸는 대신,
하나의 기구 $\mathcal{M}$에 서로 다른 제어 법칙을 적용한다.
\begin{equation}
 \vect{q}^{\star}_{k}=\pi_m(\vect{x}_k,\vect{c}_k;\mathcal{M}),
 \quad \mathcal{M}_W=\mathcal{M}_R=\mathcal{M}_C.              \label{eq:reuse-ko}
\end{equation}
평가 대상은 최대 속도 하나가 아니라
\begin{equation}
 \mathcal{O}=\{\bar v_x, e_v, C_{\rm mt},D_{\rm slip},
 P_{\rm top}(h,d),t_{\rm top},t_{\rm switch}\}                \label{eq:envelope-ko}
\end{equation}
로 나타내는 운용 범위다. 검증할 가설은 세 가지다. H1은 관절 운동과 말단 회전을
협응한 구름 주행이 동일 모델의 관절 보행 및 말단 단독 회전보다 높은 평균 전진
속도를 낸다는 것이다. H1을 명령 추종 성능이 아닌 달성 속도로 표현한 데에는
이유가 있다. 고정된 속도 명령을 초과하는 제어기는 실제로 더 빠르게 이동하면서도
추종 오차 $e_v$는 커질 수 있기 때문이다. 따라서 \ref{sec:results-ko}절에서는
평균 속도와 추종 오차를 별도로 보고한다. H2는 개방형 호 모델에서 말단 단독
회전이 실패하는 계단 형상을 관절 협응으로 통과할 수 있다는 것이다. 개구부 자체의
효과를 주장하려면 \ref{sec:method-ko}절에서 정의한 동일 조건 닫힌 바퀴 절제가
필요하다. H3은 모드 전환 중 몸체 기울기가 $60^\circ$ 한계를 넘지 않는다는 것이다.

표~\ref{tab:related-ko}의 정성적 분류 기준도 여기서 정의한다. $J_i$를 다리 $i$의
관절 속도로부터 (말단 표면을 따르는 접촉 패치의 접선 이동량, 몸체 좌표계에서 본
접촉점 위치)의 쌍으로 가는 사상이라 하자. $\operatorname{rank}J_i=1$이어서 하나의
입력이 두 값을 동시에 결정하면 \emph{결합}이라 한다. 별도의 구동기가 두 값을
독립적으로 제어하면 \emph{분리}라 한다. 본 연구처럼 세 관절이 유계 영역 안에서
두 값을 선택할 수 있지만 그중 말단 관절을 공유하는 경우는 \emph{부분 분리}로
분류한다.

\begin{table}[t]
\caption{제안 플랫폼의 형태학적 분류}
\label{tab:related-ko}
\centering\small
\resizebox{\columnwidth}{!}{%
\begin{tabular}{lcccc}
\toprule
플랫폼 & 다리 구동 & 구름 요소 & 별도 바퀴 구동 & 몸체/접점 권한\\
\midrule
RHex \cite{saranli2001rhex} & 1/다리 & 탄성 호 & 없음 & 결합\\
Whegs I \cite{quinn2002whegs} & 결합 & 3-스포크 & 공유 & 결합\\
Q-Whex \cite{zhang2023qwhex} & 1/다리 & 준바퀴 & 없음 & 결합\\
Cassino III \cite{tedeschi2017cassino} & 관절형 & 옴니휠 & 있음 & 분리\\
제안 구조 & 3/다리 & 열린 말단 호 & 없음 & 부분 분리\\
\bottomrule
\end{tabular}}
\end{table}
